% basic nastaveni
\documentclass[12pt]{article}
\setlength{\parindent}{0pt}
\setlength{\parskip}{1em}
\usepackage[utf8]{inputenc}
\usepackage[czech]{babel}
\usepackage[T1]{fontenc}
\usepackage{graphics}
\usepackage{caption}
\usepackage{hyperref}
\usepackage{xcolor}
\usepackage{float}

% matematicke package
\usepackage{amsmath}
\usepackage{amsfonts}
\usepackage{amssymb}
\usepackage{geometry}
\geometry{margin=2.5cm}
\usepackage{pgfplots}
\pgfplotsset{compat=1.18}
\usepackage{mathrsfs}

\title{Maturitní otázka číslo 17: Vektorová algebra}
\author{}
\date{}

\begin{document}
\maketitle

\subsection*{Soustavy souřadnic}
\subsubsection*{Na přímce}
Soustava souřadnic na přímce se nazývá číselná osa. Ta se zavádí pomocí přímky $p$ a dvou bodů $O$ a
$I$. Úsečka $OI$ se nazývá jednotková úsečka a její délká je jednotkou délky číselné osy. Bod $O$ se
nazývá počátek číselné osy a $I$ se nazývá jednotkový bod. Každému bodu na ose odpovídá právě jedno
reálné číslo a každému reálnému číslu odpovídá právě jeden bod na ose.
\subsubsection*{V rovině}
Soustava souřadnic v rovině se zavádí pomocí dvou os $x, y$, které jsou na sebe kolmé a mají
společný počátek v $O$. Tato soustava se nazývá kartézská soustava.
Jednotkový bod osy $x$ označíme $I$ a osy $y$ bodem $J$. Jednotkové úsečky
$OI$ a $OJ$ jsou většinou shodné. Bod $O$ se nazývá počátkem soustavy souřadnic a osy $x, y$ se
nazývají osy soustavy souřadnic. Osy $x, y$ rozdělují rovinu na 4 kvadranty. První kvadrant se nachází
vpravo nahoře od středu soustavy a pojmenování pokračuje proti směru hodinových ručiček. Nechť je bod
$A$ libovolný bod roviny. Veďme tímto bodem rovnoběžky k obou osám. Rovnoběžka s osou $y$ protíná
osu $x$ v bodě odpovídajícímu relnému číslu $a_1$. Rovnoběžka s osou $x$ protíná osu $y$ v bodě
odpovídajícímu reálnému číslu $a_2$. Číslo $a_1$ se nazývá první souřadnice nebo x-ová souřadnice
boud $A$ a $a_2$ se nazývá druhá souřadnice nebo y-ová souřadnice bodu $A$. Pokud chceme vyjadrit,
že bod $A$ ma souřadnice $a_1, a_2$, píšeme $A[a_1, a_2]$.

Kromě kartézské soustavy existuje ještě polární soustava. Pro její definici zvolíme nějakou přímku
roviny nazývanou polárni osou. Její počátek $O$ nazýváme pól. Jednotkový bod polární osy označíme
$I$. Libovolný bod $M$ roviny má polární souřadnice $O[r, \varphi]$, přičemž $r = OM$ a $\varphi$ je
polární úhle který je roven úhlu mezi kladnou poloosou $OI$ a polopřímkou $OM$. Polární úhel je
orientovaný a měří se od kartézské osy proti směru hodinových ručiček.
\subsubsection*{V prostoru}

Soustava souřadnic v prostoru se zavádí pomocí tří os $x, y$ a $z$, které jsou na sebe navzájem kolmé.
Tyto osy mají společný počátek $O$. Jednotkový bod osy $x$ značíme $I$, osy $y$ $J$ a osy $z$ $K$.
Roviny $OIJ, OJK, OKI$ se nazývají roviny souřadnic. Roviny jak označené $xy, yz, zx$ dělí prostor
souřadnic na takzvané oktanty. Když se zakresluje soustava souřadnic v prostoru, dvě nejčastejší
zobrazení jsou že x-ová souřadnice míří k nám, y-ová doprava a z-ová nahoru a nebo x-ová doprava,
y-ová nahoru a z-ová k nám. Soustavy se také dělí na pravotočivé a levotočivé. Namiřme palec směrem
kladné polosy $x$ a ukazováček kladné polosy $y$. Pokud je možné ukázat prostředníčkem na směr kladné
poloosy $z$, soustava je pravotočivá, pokud ne, soustava je levotočivá. Zápis bodu $A$ funguje
stejně jako u soustavě v rovině, ale s osou navíc. Zapisuje se jako $A[a_1, a_2, a_3]$.

\subsection*{Vzdálenost a střed dvou bodů, těžiště}
Na přímce se určuje vzdálenost dvou bodů $A[a], B[b]$ pomocí vzorce $|AB| = |b - a|$. Střed se určuje
pomocí vzorce $s = \frac{a + b}{2}$.

V rovině se určuje vzdálenost dvou bodů $A[x_1, y_1], B[x_2, y_2]$ pomocí pythagorovy věty. Vzorec
je $|AB| = \sqrt{(x_2 - x_1)^2 + (y_2 - y_1)^2}.$ Střed $S[s_1, s_2]$ se vypočítá pomocí vzorců $s_1 
= \frac{x_1 + x_2}{2}$ a $s_2 = \frac{y_1 + y_2}{2}$.

V prostoru se učuje vzdálenost dvou bodů $A[x_1, y_1, z_1], B[x_2, y_2, z_2]$ pomocí dvou složených
pythagorových vět. Vzorec je $|AB| = \sqrt{(x_2 - x_1)^2 + (y_2 - y_1)^2 + (z_2 - z_1)^2}$. Střed
$S[s_1, s_2, s_3$ se zjistí pomocí vzorců $s_1 = \frac{x_1 + x_2}{2}, s_2 = \frac{y_1 + y_2}{2}$ a 
$s_3 = \frac{z_1 + z_2}{2}$.

Vzorec těžiště pro trojúhelníky, rovnoběžníky a pravidelné n-úhelníky je $T = \frac{A_1 + A_2 + \dots
+ A_n}{n}$. Vzorec v rovině by vypadal takto $T = [t_1, t_2], t_1 = \frac{x_1 + x_2 + \dots + x_n}{n}
, t_2 = \frac{y_1 + y_2 + \dots + y_n}{n}$.
\subsection*{Vektory a operace s nimi}
\subsubsection*{Co to je vektor}
Vektor je množina všech orientovaných úseček, které mají stejnou velikost a stejný směr. Orientovaná
úsečka $AB$ je úsečka s orientací. Bod $A$ nazvěme počátečním bodem a bod $B$ nazveme koncovým bodem.
Tato úsečka se značí jako $\overrightarrow{AB}$. Z orientované úsečky s body $A[x_1, y_1],
B[x_2, y_2]$ získáme vektor $\vec{u} = (u_1, u_2), u_1 = x_2 - x_1, u_2 = y_2 - y_1$.
Vektor se dá brát jako předpis posunutí, tedy o kolik a jakým směrem se musíme posunout z bodu $A$,
abychom se dostali do bodu $B$. Velikost vektoru se získává pomocí vzorce $|\vec{u}| =
\sqrt{u_1^2 + u_2^2}$. Velikost vektoru odpovídá $|AB|$.

\subsubsection*{Sčítání, odčítání a násobení vektorů reálným číslem}
Sčítání vektorů $\vec{u} = (u_1, u_2), \vec{v} = (v_1, v_2)$ se dá představit
jako kombinaci těchto dvou posunutí. Místo toho, abychom se posunuli nejříve podle prvního vektoru
a následně druhého to zkombinujeme do jednoho posunutí, tedy $\vec{u} + \vec{v}
= (u_1 + v_1, u_2 + v_2)$.

Odčítání vektorů se dá představit jako sčítání opačného. Mějme dva vektory $\vec{u} =
(u_1, u_2), \vec{v} = (v_1, v_2)$. Opačný vektor se tvoří znegováním každé složky vektoru,
tedy opačný vektor je $\vec{v_o} = (-v_1, -v_2)$. Platí
$\vec{u} - \vec{v} = (u_1 - v_1, u_2 - v_2)$ nebo
$\vec{u} + \vec{v_o} = (u_1 - v_1, u_2 - v_2)$.

Pokud chceme vynásobit vektor reálným číslem $k$, tímto číslem si vynásobí každá složka daného
vektoru. Tedy $k \cdot \vec{u} = (k \cdot u_1, k \cdot u_2)$.

\subsubsection*{Nulový a jednotkový vektor}
Nulový vektor se většinou značí $\vec{o}$ a jeho všechny složky jsou nulové. Nulový
vektor udává, že se nepohneme nijak žádným směrem. Odečtením dvou stejných vektoruů získáváme nulový
vektor.

Jednotkový vektor je vektor, jehož velikost je 1, tedy platí $|\vec{u}| = 1$. Často 
používané jednotkové vektory jsou ty, které mají všechny složky nulové kromě jedné. Vytvoření
jednotkového vektoru z jiného vektoru, který má stjený směr funguje následovně. Pokud máme nenulový
vektor, jednotkový vektor získáme pokud ho vydělíme jeho vlastní velikostí $\hat{u} = \frac{\vec{u}}
{|\vec{u}|}$
Pokud jsme měli vektor $\vec{u} = (3, 4), |u| = \sqrt{3^2 + 4^2} = 5$ a nulový vektor bude
$\hat{u} = (\frac{3}{5}, \frac{4}{5})$. Můžeme ověřit že $0,6^2 + 0,8^2 = 1$.

\subsubsection*{Lineární kombinace a lineární závislost vektorů}
Mějme skupinu $N$ vektorů $\vec{v_1}, \vec{v_2}, \dots, \vec{v_n}$ a stejného počtu reálných čísel
$k_1, k_2, \dots, 
k_n$. Lineární kombinací těchto vektorů vznikne vektor $\vec{w} = \sum_{i = 1}^{n} k_i \cdot
\vec{v_i}$. Pomocí dvou vektorů, které nejsou nulové, stejné a rovnoběžné lze v 2D prostoru vytvořit
jakýkoliv jiný vektor. U vektoru se vlastnost toho, že vektory jsou stejné a rovnoběžné nazývá
kolineárnost. Pomocí třech vektorů, které jsou nekolineární a nekomplanární můžeme ve 3D prostoru
vytvořit pomocí lineární kombinaci vytvořit jakýkoliv vektor. Komplanárnost udává, že vektory
leží ve stejné rovině. Pro řešení

$N$ vektorů $\vec{v_1}, \vec{v_2}, \dots, \vec{v_n}$ jsou lineárně závislé, pokud existuje n-tice
realných čísel $k_1, k_2, \dots, k_n$, z nichž alespoň jedno je nenulové a platí $\sum_{i = 1}^{n}
k_i \cdot \vec{v_i} = \vec{o}$. V 2D prostou jsou vektory lineárně závislé, pokud je alespoň jeden z
nich nulový a nebo jsou dva kolineární. V 3D prostoru jsou vektory lineárně závislé, pokud je alespoň
jeden nulový, nebo jsou 2 z nich kolineární nebo jsou 3 z nich komplanární.

\subsubsection*{Skalární součin a odchylka vektorů}
Skalární součin je násobení vektorů, ve kterém je výsledek skalár(realné číslo). Mějme dva vektory
$\vec{u} = (x_1, y_1), \vec{v} = (x_2, y_2)$ platí $\vec{u} \cdot \vec{v} = |\vec{u}| \cdot |\vec{v}|
\cdot \cos \varphi$ a teké platí $\vec{u} \cdot \vec{v} = x_1 \cdot x_2 + y_1 \cdot
y_2$. Ve 3D by se jen přídalo $z_1 \cdot z_2$.
Pokud jsou oba vektory nenulové a výsledek je roven 0, vektory jsou navájem kolmé. To jde
odvodit pomocí prvního vzorce. Aby se pravá strana rovnala 0, musí být $\cos \varphi$ roven nule,
protože délky vekotů budou vždy kladná čísla. Kosínus dává nulu práve v $90^\circ$.

Důležitý trik je hledání kolmého vektoru k zadanému vektoru v 2D. Mějme vektor $\vec{u} = (x_1,
y_1)$. Vektor na něj kolmý je $\vec{v} = (-y_1, x_1)$ nebo $\vec{v} = (y_1, -x_1)$.

\subsubsection*{Vektorový součin}
Vektorový součin je násobení vektorů, které má jako součin vektor. Pro vektory $\vec{u} = (x_1, y_1,
z_1), \vec{v} = (x_2, y_2, z_2)$ platí $|\vec{w}| = |\vec{u}| \cdot |\vec{v}| \cdot \sin{\varphi}$
pokud jsou oba vektory nenulové. Vektor $\vec{w}$ je kolmý na oba vektory $\vec{u}$ a $\vec{v}$.
Predpis vektoru získámé pomocí determinantu matice. Pro vektory $\vec{u} = (x_1, y_1, z_1), \vec{v} =
(x_2, y_2, z_2)$ platí $\vec{w} = \vec{u} \times \vec{v} = (y_1 z_2 i - z_1 y_2 i, z_1 x_2 j - x_1
z_2 j, x_1 y_2 k - y_1 x_2 k)$.
\begin{center}
\begin{vmatrix}
    i & j & k \\
    x_1 & y_1 & z_1 \\
    x_2 & y_2 & z_2
\end{vmatrix}
\end{center}

Vektory $i, j, k$ jsou jednotkové vektory ve směru os $x, y, z$. Po přechodu do souřadnicového systému
dostaneme $\vec{w} = (y_1 z_2 - z_1 y_2, z_1 x_2 - x_1 z_2, x_1 y_2 - y_1 x_2)$. Jde si to představit,
jako odčítání dvou diagonaál od sebe znázorněných na těchto maticích.
\begin{center}
\begin{vmatrix}
    \color{purple}{i} & j & k \\
    x_1 & \color{red}{y_1} & \color{blue}{z_1} \\
    x_2 & \color{blue}{y_2} & \color{red}{z_2}
\end{vmatrix}
\begin{vmatrix}
    i & \color{purple}{j} & k \\
    \color{blue}{x_1} & y_1 & \color{red}{z_1} \\
    \color{red}{x_2} & y_2 & \color{blue}{z_2}
\end{vmatrix}
\begin{vmatrix}
    i & j & \color{purple}{k} \\
    \color{red}{x_1} & \color{blue}{y_1} & z_1 \\
    \color{blue}{x_2} & \color{red}{y_2} & z_2
\end{vmatrix}
\end{center}

Pomocí vektorového součinu lze vypočítat obsah a objem různých těles. Smíšený součin je součin, ve
kterém je jak vektorový tak skalární součin.

Rovnoběžník: $S = |\vec{u} \times \vec{v}|$

Trojúhelník: $S = \frac{1}{2}|\vec{u} \times \vec{v}|$

Rovnoběžnostěn: $V = |(\vec{u} \times \vec{v}) \cdot \vec{w}|$

Čtyřstěn: $V = \frac{1}{6}|(\vec{u} \times \vec{v}) \cdot \vec{w}|$

Hranol s trojúhelníkovou podstavou: $V = \frac{1}{2}|(\vec{u} \times \vec{v}) \cdot \vec{w}|$
\end{document}
